\documentclass[12pt,a4paper]{article}

% ============================================================
% PACKAGES
% ============================================================
\usepackage[utf8]{inputenc}
\usepackage[T1]{fontenc}
\usepackage[english]{babel}
\usepackage{geometry}
\usepackage{graphicx}
\usepackage{hyperref}
\usepackage{booktabs}
\usepackage{tabularx}
\usepackage{longtable}
\usepackage{listings}
\usepackage{xcolor}
\usepackage{amsmath}
\usepackage{setspace}
\usepackage{indentfirst}
\usepackage[numbers]{natbib}

% ============================================================
% PAGE SETUP
% ============================================================
\geometry{left=3cm, right=2cm, top=3cm, bottom=2cm}
\onehalfspacing

% ============================================================
% CODE LISTING STYLE
% ============================================================
\definecolor{codebg}{rgb}{0.95,0.95,0.95}
\definecolor{codegreen}{rgb}{0.0,0.5,0.0}
\definecolor{codegray}{rgb}{0.5,0.5,0.5}
\definecolor{codeblue}{rgb}{0.0,0.0,0.7}

\lstset{
    backgroundcolor=\color{codebg},
    basicstyle=\ttfamily\small,
    breaklines=true,
    captionpos=b,
    commentstyle=\color{codegreen},
    keywordstyle=\color{codeblue},
    numberstyle=\tiny\color{codegray},
    stringstyle=\color{red},
    frame=single,
    numbers=left,
    numbersep=5pt,
    showstringspaces=false,
    tabsize=2,
}

% ============================================================
% TITLE
% ============================================================
\title{
    \textbf{GraphRAG for Conversational History Class Data:}\\
    \large A Custom Knowledge Graph and Retrieval Architecture\\with Multi-Pass Extraction
}
\author{Eduardo Birkheuer\\
    \small Universidade Federal do Rio Grande do Sul (UFRGS)\\
    \small Instituto de Inform\'{a}tica --- Ci\^{e}ncia da Computa\c{c}\~{a}o
}
\date{2026}

% ============================================================
\begin{document}
\maketitle

\begin{abstract}
This work presents the design and implementation of a custom GraphRAG system for transforming English history class transcriptions into a queryable knowledge graph. The system extracts structured knowledge---entities, relationships, claims, and predictions---from approximately 134 episodes of lecture content using a multi-pass extraction pipeline with schema-constrained structured outputs. Extracted knowledge is stored in a Neo4j property graph with bi-temporal metadata, and a hybrid retrieval layer combines vector similarity search, graph traversal, text-to-Cypher generation, and community-based thematic summaries, exposed to client-side LLMs via an MCP server. The architecture is custom-built rather than framework-dependent, justified by three empirical findings: schema-constrained extraction achieves 98.8\% precision versus open extraction's 40\% predicate error rate; the LoCoMo benchmark demonstrates that transforming dialogue into structured assertions dramatically outperforms raw transcript retrieval; and no existing framework natively supports the intersection of claim reification, bi-temporal modeling, and domain-specific historical ontology that this system requires.
\end{abstract}

\tableofcontents
\newpage

% ============================================================
% Include sections
% ============================================================
\section{Background and Technologies}
\label{sec:conceitos}

This section presents the theoretical concepts and technologies required to understand the proposed architecture.

\subsection{Retrieval-Augmented Generation (RAG)}

Large Language Models (LLMs) have a fundamental limitation: their knowledge is static, bounded by training data. Additionally, LLMs frequently produce factually incorrect responses---a phenomenon known as hallucination~\cite{huang2023survey}. Retrieval-Augmented Generation (RAG) addresses both problems by combining an information retrieval component with the generative model~\cite{lewis2020rag}.

The typical RAG pipeline operates in three stages:

\begin{enumerate}
    \item \textbf{Indexing:} Documents are split into chunks ($\sim$500--1500 tokens) and converted into high-dimensional numerical vectors (embeddings) via models such as OpenAI's \texttt{text-embedding-3-small}. These vectors are stored in a vector database.
    
    \item \textbf{Retrieval:} When a user submits a query, it is also converted into an embedding. The system searches the vector store for chunks whose embeddings are most similar to the query (typically via cosine similarity).
    
    \item \textbf{Generation:} Retrieved chunks are injected into the LLM's prompt as context, and the model generates a response grounded in this information.
\end{enumerate}

RAG represents a significant advance: it enables access to current information without retraining, grounds responses in verifiable sources, and reduces hallucinations. However, the vector-based approach has structural limitations. Text chunks are treated as flat, independent units---all relational information between entities, causality between events, and argumentative structure is lost in the fragmentation process. For domains where relationships and structure are central to understanding---such as history, where events cause other events and figures interact in complex networks---these limitations become especially problematic.

\subsection{Knowledge Graphs}

A Knowledge Graph (KG) is a data structure that represents information as a graph, where nodes represent entities (people, places, concepts) and edges represent relationships between them~\cite{hogan2021knowledge}. The basic unit of information is the triple: (subject)--[predicate]$\rightarrow$(object). For example: (Alexander the Great)--[CONQUERED]$\rightarrow$(Persian Empire).

Two primary paradigms exist for knowledge graph representation:

\textbf{RDF (Resource Description Framework):} A W3C standard based on subject-predicate-object triples. Uses URIs for global identification, formal ontologies (OWL) for logical reasoning, and SPARQL as query language. Ideal for interoperability and Linked Data, but requires complex mechanisms (blank nodes, reification via named graphs) for attaching properties to relationships.

\textbf{Property Graphs:} A model where both nodes and edges can carry arbitrary properties (key-value pairs). The predominant query language is Cypher, based on visual pattern matching. Cypher has been internationally standardized as GQL (ISO/IEC 39075:2024)---the first new ISO database language since SQL in 1987.

This work adopts the property graph model for three reasons: (1) schema flexibility---properties can be added to nodes and edges without restructuring; (2) query expressiveness---Cypher is more natural for multi-hop traversal than SPARQL; (3) alignment with the Neo4j ecosystem, which offers native vector indexes, community detection via GDS, and mature LLM integration.

\textbf{Ontology} refers to the schema defining which entity types and relationship types are permitted in the graph. A well-defined ontology serves as a contract: it guides extraction (the LLM knows what to look for), prevents inconsistencies (invalid types are rejected), and enables predictable queries (consumers know what exists in the graph).

\textbf{Reification} is the pattern of transforming a relationship into an intermediate node to allow it to have its own properties and relationships. Instead of modeling a claim as a property on an edge between two entities, a Claim node is created that connects to its source utterance, the episode in which it was stated, and the entities it references. Neo4j documents this pattern as the ``Intermediate Node Pattern'' in its official data modeling course.

\subsection{GraphRAG}

GraphRAG combines Knowledge Graphs with Retrieval-Augmented Generation, using graph structure to enrich the context provided to the LLM beyond what pure vector search offers.

The seminal Microsoft work~\cite{edge2024graphrag} established the term and demonstrated its efficacy in \textit{sensemaking} tasks---global comprehension of large corpora. The Microsoft GraphRAG architecture operates in two phases:

\textbf{Indexing:} Texts are processed by an LLM that extracts entities and relationships, building a knowledge graph. Over this graph, the Leiden algorithm detects communities (clusters of densely connected entities). Each community receives an LLM-generated summary at multiple hierarchy levels.

\textbf{Querying:} ``Local'' queries (about specific entities) use entity search with neighborhood expansion. ``Global'' queries (about general themes) synthesize summaries of relevant communities.

GraphRAG approaches outperform vector-based RAG in 70--80\% of complex sensemaking tasks~\cite{edge2024graphrag}. However, the Microsoft implementation has limitations: it recognizes only 4 entity types (Person, Organization, Geography, Event); consumes $\sim$610,000 tokens per global query; requires full reindexing for updates; and has no optimizations for conversational data.

Relevant alternatives include \textbf{LightRAG}~\cite{guo2024lightrag}, which uses open extraction with $\sim$100 tokens per query (6,000$\times$ more efficient), and \textbf{LazyGraphRAG}~\cite{edge2024lazygraphrag}, which eliminates upfront summarization at 0.1\% of full GraphRAG indexing cost, achieving 95/96 statistical wins against all alternatives.

The concept of \textbf{hybrid retrieval}---combining vector search (semantic) with graph traversal (structural)---is central to this work. Vector search finds relevant entry points; graph traversal expands context with relationships, claims, and connected entities that pure semantic similarity would not capture.

\subsection{Neo4j and Cypher}

Neo4j is the most widely used graph database globally, implementing the property graph model. Its relevance to this work stems from specific capabilities:

\textbf{Cypher} is Neo4j's query language, based on visual pattern matching where graph patterns are expressed with ASCII art: \texttt{(a:Person)-[:KNOWS]->(b:Person)}. Version 5.x introduced native vector query support, eliminating the need for a separate vector database.

\textbf{Vector Indexes} enable semantic similarity search directly within Neo4j. Each node can store an embedding vector and be queried via approximate nearest neighbor (ANN):

\begin{lstlisting}[language=SQL, caption={Neo4j vector similarity query}]
CALL db.index.vector.queryNodes(
  'utterance_embeddings', 5, $queryEmbedding
) YIELD node AS chunk, score
\end{lstlisting}

\textbf{MERGE} is the idempotent upsert operator: it creates a node/relationship if it does not exist, or finds the existing one if already created. With \texttt{ON CREATE SET} and \texttt{ON MATCH SET} clauses, it enables progressive enrichment without duplication---essential for an incremental ingestion pipeline.

\textbf{Neo4j GDS (Graph Data Science)} is a library of graph algorithms including community detection (Leiden, Louvain), centrality (PageRank, Betweenness), and similarity. For this work, the Leiden algorithm detects thematic communities in the entity graph.

\subsection{Knowledge Extraction with LLMs}

The construction of Knowledge Graphs from unstructured text traditionally relied on NLP pipelines (NER + Relation Extraction + Coreference Resolution). With LLMs, end-to-end extraction via prompting became possible---feeding text and instructions, the model returns structured entities and relationships.

\textbf{Single-pass vs.\ Multi-pass.} Requesting that the LLM extract entities, relationships, claims, and resolve coreferences in a single call produces suboptimal results---the model does everything superficially. Research converges on multi-pass superiority: ECR demonstrates +16\% in triple F1 (single vs.\ multi-pass); KGGen~\cite{mo2025kggen} achieves 18\% improvement over GraphRAG via multi-stage extraction (extract$\rightarrow$cluster$\rightarrow$normalize); SAC-KG~\cite{chen2024sackg} reaches 89.32\% precision with a Generator$\rightarrow$Verifier$\rightarrow$Pruner pipeline at 1M+ node scale.

\textbf{Structured Outputs} are a constrained decoding mechanism offered by LLM providers (OpenAI Structured Outputs, August 2024) that guarantees 100\% compliance with a provided JSON Schema. Unlike asking the model to ``return JSON'' (compliance $<$40\% with GPT-4-0613), structured outputs use restrictive grammars during decoding, making it impossible to generate output outside the schema. Benchmarks (JSONSchemaBench, ICML 2025) demonstrate that constrained decoding not only guarantees format but improves downstream quality by up to 4\%.

\textbf{Entity Resolution} is the process of identifying when different names refer to the same real-world entity. Modern approaches use cascades from cheap to expensive: text normalization (40\% of duplicates), fuzzy string matching via \texttt{token\_sort\_ratio} (35\%), embedding similarity (20\%), and LLM judgment for only the remaining 5\% of genuinely ambiguous pairs. CORE-KG demonstrated that omitting coreference resolution causes a 28.32\% increase in node duplication.

\subsection{Community Detection and the Leiden Algorithm}

Community detection in graphs identifies sets of nodes that are more densely connected to each other than to the rest of the graph. Intuitively, communities correspond to thematic clusters---in a history graph, a community might group ``Roman Empire,'' ``Julius Caesar,'' ``Roman Senate,'' and ``Roman Civil War.''

The \textbf{Louvain algorithm}~\cite{blondel2008louvain} is the classic method: it optimizes modularity by iteratively moving nodes between communities. However, Louvain can produce poorly-connected communities (internally disconnected).

The \textbf{Leiden algorithm}~\cite{traag2019leiden} resolves this guarantee: beyond optimizing modularity, it ensures all detected communities are internally well-connected. Leiden adds a refinement phase that Louvain lacks, resulting in superior partition quality.

\textbf{Hierarchical communities} are obtained by running Leiden with intermediate communities enabled, producing partitions at multiple resolution levels: Level 0 (granular---``Alexander's campaigns in Persia''), Level 1 (medium---``Hellenistic Conquests''), Level 2 (broad---``Classical Antiquity''). The gamma parameter controls resolution: larger values produce smaller, more numerous communities.

For \textbf{incremental updates}, the seeded Leiden mechanism reuses the previous partition as a starting point. When a new episode adds entities to the graph, only affected communities are repartitioned---avoiding recomputing the entire graph. Dynamic Frontier Leiden~\cite{sahu2024dfleiden} formalizes this process, achieving comparable modularity to static Leiden at dramatically reduced runtime.

\subsection{Model Context Protocol (MCP)}

The Model Context Protocol (MCP) is an open standard for connecting AI systems to external data sources. Announced by Anthropic in November 2024 and subsequently donated to the Agentic AI Foundation under the Linux Foundation (December 2025, co-founded with Block and OpenAI), the protocol reached 97 million monthly SDK downloads and 10,000+ servers by late 2025.

MCP operates over JSON-RPC 2.0 and defines three primitives:

\begin{itemize}
    \item \textbf{Tools:} Functions the LLM can invoke (e.g., search the graph, compare entities)
    \item \textbf{Resources:} Static data the LLM can read (e.g., the graph schema)
    \item \textbf{Prompts:} Reusable templates for common interactions
\end{itemize}

The advantage of MCP over traditional REST APIs is that the client-side LLM can autonomously decide which tool to invoke based on the user's question. The system becomes self-navigable without a dedicated UI---the LLM \textit{is} the interface.

\section{Related Work}
\label{sec:related}

This section positions the present work in the context of existing literature, identifying the gaps that motivate the construction of a custom pipeline.

\subsection{Existing GraphRAG Systems}

\textbf{Microsoft GraphRAG}~\cite{edge2024graphrag} is the seminal work that formalized the combination of knowledge graphs with RAG for sensemaking tasks. The system extracts entities and relationships via LLM, applies hierarchical Leiden community detection, and generates per-community summaries. Two query modes---local (entity-centric) and global (community summaries)---cover point and thematic queries respectively.

Despite its conceptual contribution, Microsoft GraphRAG presents significant practical limitations. Its default extraction recognizes only four entity types (Person, Organization, Geography, Event), insufficient for capturing stances, predictions, and causal analyses present in history lectures. The global query mode consumes approximately 610,000 tokens per query---versus $\sim$100 tokens for LightRAG, a 6,000$\times$ difference. The official repository warns: ``Using GraphRAG with your data out of the box may not yield the best possible results.'' Shereshevsky (2025) reports indexing costs of \$33K for large datasets, infeasible for a 134-episode corpus requiring incremental updates.

\textbf{LightRAG}~\cite{guo2024lightrag} proposes a lightweight alternative: open extraction without schema constraints, with incremental updates via graph union and efficient queries ($\sim$100 tokens). Its simplicity comes at a cost: zero schema constraints produce inconsistent relationship types and entity duplication. The Plumber framework analysis demonstrates that open extraction results in 40\% predicate errors. The paper was withdrawn from ICLR 2025 review.

\textbf{LazyGraphRAG}~\cite{edge2024lazygraphrag} eliminates upfront summarization, with indexing cost at 0.1\% of full GraphRAG. In evaluation, it achieved 95 of 96 statistical wins against Vector RAG, GraphRAG Local/Global/DRIFT, LightRAG, RAPTOR, and TREX. It is the default recommendation for non-conversational document corpora. However, it possesses no features for dialogue data: no speaker attribution, no temporal modeling, no claim reification.

\textbf{GraphRAG-Bench}~\cite{xiang2026graphragbench} provides the first rigorous cross-framework benchmark, concluding that GraphRAG frequently underperforms vanilla RAG on real-world tasks, offering guidelines on when graph structure provides measurable benefits---typically in multi-hop, comparative, and global synthesis queries.

\subsection{Conversational Memory Systems}

\textbf{Graphiti/Zep}~\cite{rasmussen2025graphiti} is the most mature system for conversational knowledge graphs. With 22,800+ GitHub stars and 25,000+ weekly PyPI downloads, it defines the graph as $G = (N, E, \varphi)$ with three subgraphs: Episode (preserves raw data), Entity (extracted semantic knowledge), and Community (thematic summaries).

Graphiti's central innovation is its bi-temporal model: every semantic edge carries four timestamps---\texttt{created\_at} (when ingested), \texttt{expired\_at} (when invalidated in system), \texttt{valid\_at} (when fact became true in reality), and \texttt{invalid\_at} (when fact stopped being true). When contradictory information arrives, the old edge is invalidated and the fact text is regenerated to reflect evolution.

Despite its temporal sophistication, Graphiti is designed for personal assistant memory---tracking user preferences, personal facts, previous conversations. Its entity extraction is generic, without domain ontology. For specific historical analysis (figures, empires, battles, causal relationships between events), fine-grained control over entity and relationship types is required.

\textbf{SGMem}~\cite{wu2025sgmem} represents dialogue as sentence-level graphs within chunked units, capturing turn-, round-, and session-level associations, outperforming baselines on LongMemEval and LoCoMo.

\textbf{MAGMA}~\cite{jiang2026magma} extends the paradigm with orthogonal graphs (semantic, temporal, causal, entity) per memory item, achieving 45.5\% higher reasoning accuracy via policy-guided traversal.

\textbf{``Does Memory Need Graphs?''}~\cite{hu2026memorygraphs} introduces a unified framework decomposing dialog memory systems into core components. Its key finding: sentences-as-nodes can outperform entity-level KG approaches for conversational memory.

\subsection{Knowledge Extraction from Dialogue}

\textbf{OKE-D}~\cite{rook2025oked} (Rook, ISWC 2025 Doctoral Consortium) is the most directly relevant academic reference. It formalizes Open Knowledge Extraction from Dialogue and identifies four unique structural challenges: (1) cross-turn assertions with higher subject-object span distances than documents; (2) anaphora and pronoun density requiring robust coreference resolution; (3) dialogue acts where not all utterances assert facts; (4) multiple perspectives where speakers express conflicting views.

Rook proposes a perspective-aware dialogue ontology integrating speaker identity, temporal context, and dialogue acts via SIO and PROV-O. The work demonstrates that even frontier models (Gemini 2.5 Pro) fail to extract important knowledge from dialogues without specialized prompting.

\textbf{LoCoMo}~\cite{maharana2024locomo} is the most relevant long-term conversational memory benchmark, evaluating over $\sim$300 turns and up to 35 sessions. Key findings: human F1 = 87.9\%, GPT-3.5-16K = 37.8\%, best RAG = 41.4\%---LLMs lag 56\% behind humans. The most critical finding: ``RAG does particularly well when dialogues are transformed into a database of assertions about each speaker's life and persona''---direct validation of the knowledge graph extraction approach.

\textbf{KGGen}~\cite{mo2025kggen} demonstrates multi-stage extraction (extract$\rightarrow$cluster$\rightarrow$normalize) scoring 66.07\% on the MINE benchmark versus GraphRAG's 47.80\% and OpenIE's 29.84\%---an 18\% improvement.

\textbf{SAC-KG}~\cite{chen2024sackg} proposes a Generator$\rightarrow$Verifier$\rightarrow$Pruner pipeline achieving 89.32\% precision at 1M+ node scale, representing $>$20\% increase over prior state of the art.

\textbf{CORE-KG} demonstrates via ablation that removing coreference resolution causes a 28.32\% increase in node duplication, and that removing structured prompts causes 73\% more noise---quantitative evidence for each pipeline component.

\subsection{Claim Reification and Argumentation Frameworks}

The modeling of claims as first-class nodes in knowledge graphs is grounded in multiple traditions:

\textbf{AIF (Argument Interchange Format)} defines I-Nodes (claims/information), S-Nodes (scheme applications subdivided into inference, conflict, and preference), and L-Nodes (speech acts in dialogue). An AIF argument graph is a bipartite-like structure where edges always connect I-Nodes to S-Nodes and vice versa.

\textbf{DialAM-2024}~\cite{ruizdolz2024dialam} presents the first shared task in dialogical argument mining using Inference Anchoring Theory, modeling dialogue with locution nodes (what was said), proposition nodes (claims), and relation nodes (argumentative/illocutionary).

\textbf{Discourse Graphs}~\cite{chan2024discourse} model knowledge claims linked to supporting evidence through typed edges (supports, opposes, informs), with 3 years of practical deployment in distributed research communities.

\textbf{Kialo} demonstrates practical scalability: 18,000+ debates and 720,000+ claims in production, using a tree structure where each claim is Pro or Con relative to its parent.

The \textbf{Neo4j Intermediate Node Pattern} provides the technical basis for reification in property graphs: replacing a relationship that needs rich metadata with a node, enabling arbitrary properties, relationships to other claims, and independent lifecycle tracking.

\subsection{Comparative Analysis}

Table~\ref{tab:comparison} compares the proposed system with existing alternatives across dimensions relevant to the history lecture use case.

\begin{table}[htbp]
\centering
\caption{Comparative analysis of GraphRAG systems}
\label{tab:comparison}
\small
\begin{tabularx}{\textwidth}{lXXXX}
\toprule
\textbf{Dimension} & \textbf{This work} & \textbf{MS GraphRAG} & \textbf{LightRAG} & \textbf{Graphiti} \\
\midrule
Claim reification & Native (first-class nodes) & Not supported & Not supported & Not supported \\
Temporal modeling & Bi-temporal & None & None & 4 timestamps \\
Schema control & Full (custom ontology) & 4 fixed types & Open extraction & Generic entities \\
Extraction quality & Schema-constrained & Generic & 40\% errors & Generic \\
Cost (134 episodes) & \$15--45 & \$33K (est.) & \$15--30 & Variable \\
Incremental updates & Idempotent MERGE & Partial reindex & Graph union & Dynamic \\
Community detection & Leiden (hierarchical) & Leiden & None & Label prop. \\
Tokens per query & $\sim$100--1000 & $\sim$610,000 & $\sim$100 & Variable \\
\bottomrule
\end{tabularx}
\end{table}

\subsection{Positioning}

No existing system addresses the intersection of: (1) claim modeling as first-class entities with claim-to-claim relationships; (2) multi-pass extraction with schema constraints for the historical domain; (3) temporal tracking of verifiable geopolitical predictions; (4) adaptive hybrid retrieval by query type. The present work fills this gap by constructing a custom pipeline incorporating the best patterns from each system---adapted for a history lecture domain where claims, predictions, and causal relationships constitute the highest-value analytical content.

\section{System Design and Implementation}
\label{sec:development}

This section describes the architecture, implementation, and technical decisions of the proposed system.

\subsection{Architecture Overview}

The system is organized in four layers with well-defined responsibilities: a Client Layer (any MCP-compatible LLM), an MCP Server (6 domain tools), a Retrieval Layer (4 adaptive strategies), and a Neo4j Graph (vector + structural indexes), fed by an Ingestion Pipeline (download $\rightarrow$ transcribe $\rightarrow$ clean $\rightarrow$ chunk $\rightarrow$ extract $\rightarrow$ resolve $\rightarrow$ upsert $\rightarrow$ embed $\rightarrow$ communities).

The decision to build a custom pipeline (versus adopting a framework) is justified by converging evidence: (1) schema-constrained extraction achieves 98.8\% precision versus 40\% error with open extraction (ODKE+, Plumber); (2) no existing framework models claims as first-class nodes; (3) the estimated cost of \$15--45 for 134 episodes makes the project academically viable versus \$33K for Microsoft GraphRAG at similar scale.

The technology stack consists of: Python 3.11+ for all logic; OpenAI API exclusively for LLM calls and embeddings (eliminating variability between providers); Neo4j 5.x as graph database (native vector indexes, GDS for Leiden, Cypher for queries); FastMCP for the MCP server.

\subsection{Knowledge Graph Schema}

\subsubsection{Design Principles}

The schema follows three principles derived from the literature:

\textbf{Conversation-centric hierarchy.} The atomic unit is a lecture chunk, not a document paragraph. The hierarchy Episode $\rightarrow$ Segment $\rightarrow$ Utterance mirrors Graphiti's Episode subgraph~\cite{rasmussen2025graphiti} and Di Bratto et al.'s (2021) utterance chaining pattern, preserving conversational structure that flat chunking destroys.

\textbf{Claims as first-class nodes (reification).} Claims, predictions, and opinions are modeled as intermediate nodes following Neo4j's official Intermediate Node Pattern, AIF's I-Node formalism, and Kialo's production-validated claim-tree architecture. Reification enables: arbitrary metadata (confidence, stance, verification status), claim-to-claim relationships (supports, contradicts, revises), and traceable provenance to the exact lecture moment.

\textbf{Bi-temporal metadata.} Every claim carries two temporal dimensions: when it was stated (via Episode publishDate, through the content hierarchy) and what time period it references (via explicit date properties on claims/predictions).

\subsubsection{Node Types}

The ontology defines 9 node types in 3 layers:

\textbf{Content Layer:} Episode (a single lecture), Segment (coherent topic block, $\sim$10--15 minutes), Utterance (text chunk, $\sim$1500 tokens---the vector search unit).

\textbf{Entity Layer:} Person (historical figures, leaders, philosophers), Organization (empires, nations, dynasties, religions), Topic (concepts, ideologies, geographic regions), Event (wars, battles, treaties, revolutions).

\textbf{Analytical Layer:} Claim (historical interpretation, geopolitical prediction, or causal assertion by Prof.\ Jiang), Community (thematic cluster detected by Leiden).

\subsubsection{Relationship Types}

Key relationships include:
\begin{itemize}
    \item Content hierarchy: \texttt{HAS\_SEGMENT}, \texttt{HAS\_CHUNK}, \texttt{NEXT\_SEGMENT}, \texttt{NEXT\_CHUNK}
    \item Historical dynamics: \texttt{ROLE\_IN}, \texttt{PARTICIPATED\_IN}, \texttt{LED\_TO}, \texttt{CAUSED}
    \item Entity grounding: \texttt{MENTIONS} (Utterance $\rightarrow$ Entity)
    \item Claim network: \texttt{EXTRACTED\_FROM}, \texttt{STATED\_DURING}, \texttt{ABOUT}, \texttt{SUPPORTS}, \texttt{CONTRADICTS}, \texttt{REVISES\_PREDICTION}
    \item Communities: \texttt{IN\_COMMUNITY}, \texttt{PARENT\_COMMUNITY}
\end{itemize}

\subsubsection{Expected Scale}

For 134 episodes: $\sim$134 Episode nodes, $\sim$1,300 Segments, $\sim$30,000 Utterances, $\sim$100--300 Persons, $\sim$100 Organizations, $\sim$100--200 Topics, $\sim$200 Events, $\sim$3,000--8,000 Claims, and $\sim$80,000--250,000 relationships---comfortably within Neo4j Community Edition capacity.

\subsection{Preprocessing Pipeline}

\subsubsection{Download and Transcription}

Audio is extracted from YouTube videos via \texttt{yt-dlp} (mp3 format, 64kbps---sufficient for speech) and transcribed using Groq's Whisper large-v3. For $\sim$1-hour lectures exceeding Groq's 25~MB limit, audio is segmented into 20-minute blocks via \texttt{ffmpeg}, transcribed individually, and concatenated. The pipeline is idempotent: existing transcriptions are skipped.

\subsubsection{Text Cleaning}

Raw audio transcriptions contain artifacts that waste tokens and confuse extraction: (1) Unicode normalization; (2) removal of transcription artifacts (\texttt{[inaudible]}, timestamps); (3) removal of filler words (``uh,'' ``you know''); (4) stutter removal (``the the Roman'' $\rightarrow$ ``the Roman''); (5) whitespace normalization. The cleaning is conservative: only obvious noise is removed without altering semantic content.

\subsubsection{Two-Pass Chunking}

Fixed-size chunking destroys segment boundaries. The system employs a two-pass strategy:

\textbf{Level 1---LLM Segmentation:} Text is split into sentences (via regex with abbreviation protection). Sentences are grouped into sliding windows of 20 with overlap of 3. For each window, GPT-4o-mini identifies topic boundaries. Boundaries are deduplicated and used to construct Segments---variable-length blocks, each about one coherent topic.

\textbf{Level 2---Deterministic Token Split:} Each Segment is split into chunks of at most 1500 tokens. Splitting never occurs mid-sentence. Three sentences overlap between consecutive chunks, ensuring cross-boundary assertions are captured. Chunks below 200 tokens are merged into the previous chunk.

The 1500-token target balances extraction quality (more context per call) against cost (fewer tokens per call). The 3-sentence overlap is based on Rook's (ISWC 2025) finding about higher subject-object span distances in dialogue data.

\subsection{Multi-Pass Extraction Pipeline}

\subsubsection{Justification}

Single-pass extraction is documented to be inferior. Research converges: ECR demonstrates +16\% in triple F1; KGGen~\cite{mo2025kggen} +18\% over GraphRAG; SAC-KG~\cite{chen2024sackg} 89.32\% precision with three components. Each pass with a focused task allows the LLM to concentrate its capacity on one problem at a time.

\subsubsection{Pass Architecture}

The pipeline uses five sequential passes with tiered model selection, as shown in Table~\ref{tab:passes}.

\begin{table}[htbp]
\centering
\caption{Multi-pass extraction architecture}
\label{tab:passes}
\begin{tabular}{clll}
\toprule
\textbf{Pass} & \textbf{Task} & \textbf{Model} & \textbf{Justification} \\
\midrule
1 & Entity extraction & GPT-4o-mini & High volume, cheap \\
2 & Relationship extraction & GPT-4o-mini & Same chunks, entities as context \\
3 & Claim/opinion extraction & GPT-4o & Nuanced stance detection \\
4 & Coreference resolution & GPT-4o & Deep contextual reasoning \\
5 & Entity resolution & Embeddings + GPT-4o-mini & Only 5\% need LLM \\
\bottomrule
\end{tabular}
\end{table}

\subsubsection{Pass Details}

\textbf{Pass 1 (Entity Extraction):} Input is a chunk plus episode context. Output is a list of entities with canonical name, type (constrained to person/organization/topic/event), description, aliases, and type-specific properties. The schema is enforced via OpenAI Structured Outputs---the model is physically unable to return an invalid entity type.

\textbf{Pass 2 (Relationship Extraction):} Same chunk plus entities from Pass 1 injected as context. The model focuses on finding relationships between \textit{known} entities rather than re-extracting them. Relationship types are constrained to 6 schema-valid types.

\textbf{Pass 3 (Claim Extraction):} The most critical pass---identifies what Prof.\ Jiang argues, interprets, predicts, and analyzes. Distinguishes factual assertions, historical interpretations, geopolitical predictions, and causal analyses. GPT-4o is used because stance detection requires nuanced understanding.

\textbf{Pass 4 (Coreference Resolution):} Operates on a sliding window of 4 consecutive chunks, resolving ambiguous references (``he,'' ``the empire'') to canonical entity names from Pass 1. The window is necessary because pronouns often refer to entities in previous chunks.

\subsubsection{Structured Outputs and Caching}

All passes use OpenAI Structured Outputs with Pydantic schemas: responses are guaranteed valid JSON conforming to the defined schema. Every LLM call is cached by \texttt{hash(system\_prompt + user\_prompt + model)}: reprocessing costs \$0 (cache hit), enabling idempotent pipelines and facilitating debugging via direct inspection of cached responses.

\subsection{Entity Resolution}

After extraction, the same entity appears with different names across chunks and episodes. The resolution module uses a cascade routing from cheap to expensive:

\begin{enumerate}
    \item \textbf{Exact match} ($\sim$40\% of duplicates): Normalizes case, whitespace, accents. Cost: zero.
    \item \textbf{Fuzzy string matching} ($\sim$35\%): RapidFuzz \texttt{token\_sort\_ratio} $>$ 85 between same-type entities. Captures ``Philip the Second'' $\approx$ ``Philip II.'' Cost: negligible.
    \item \textbf{Embedding similarity} ($\sim$20\%): Cosine similarity $>$ 0.85 between embeddings of ``name: description.'' Captures ``USA'' $\approx$ ``United States'' when descriptions overlap. Cost: one embedding call per entity.
    \item \textbf{LLM judgment} ($\sim$5\%): Only genuinely ambiguous pairs go to GPT-4o-mini for merge/no-merge decision with justification.
\end{enumerate}

The guiding principle is conservative: \textbf{when in doubt, do not merge}. A false positive (incorrect merge) corrupts the graph; a false negative (keeping duplicates) is harmless and correctable. The output is a merge map with transitive closure resolved.

\subsection{Graph Upsert}

Insertion follows Neo4j's MERGE pattern with the fundamental rule: never MERGE entire patterns with unbound variables; always bind nodes first, then MERGE relationships between bound nodes.

Insertion order ensures endpoints exist before creating relationships: (1) Episode, (2) Segments + chains, (3) Utterances + chains, (4) Entities + MENTIONS, (5) Claims + provenance/grounding, (6) Entity-Entity relationships.

IDs are deterministic: \texttt{make\_entity\_id("person", "Alexander the Great")} $\rightarrow$ \texttt{"person-alexander-the-great"}. Same input always produces the same ID, ensuring MERGE finds existing nodes. Progressive enrichment via \texttt{ON MATCH SET} updates descriptions (keeping the longer one) and accumulates aliases.

\subsection{Community Detection and Summarization}

After all entities and relationships are inserted, the Leiden algorithm is executed via Neo4j GDS:

\begin{enumerate}
    \item \textbf{Projection:} Entity graph (Person, Organization, Topic) with co-occurrence relationships projected in-memory with UNDIRECTED orientation.
    \item \textbf{Leiden:} Executed with hierarchical output and fixed random seed for reproducibility.
    \item \textbf{Summarization:} For each community, members and internal relationships are collected and sent to an LLM that generates a descriptive thematic summary.
    \item \textbf{Storage:} Summaries saved as Community nodes with embeddings for vector retrieval.
\end{enumerate}

For incremental updates, seeded Leiden reuses the previous partition. Only communities whose membership changed need re-summarization---typically 5--15 communities per new episode.

\subsection{Retrieval Layer}

\subsubsection{Query Classification}

A lightweight classifier (GPT-4o-mini, temperature 0, structured output) categorizes each query into one of four types, each routed to an optimal strategy: factual (vector + graph fan-out), comparative (Cypher template by entity), temporal (chronological claim aggregation), or thematic (community summaries).

\subsubsection{Vector + Graph Fan-Out}

The core retrieval pattern: vector search finds semantically relevant chunks; Cypher traversal expands context with episode metadata, connected entities, and extracted claims:

\begin{lstlisting}[language=SQL, caption={Vector search with graph fan-out}]
CALL db.index.vector.queryNodes(
  'utterance_embeddings', $topK, $queryEmbedding
) YIELD node AS utterance, score
OPTIONAL MATCH (utterance)<-[:HAS_CHUNK]-(seg)<-[:HAS_SEGMENT]-(ep)
OPTIONAL MATCH (claim:Claim)-[:EXTRACTED_FROM]->(utterance)
OPTIONAL MATCH (claim)-[:ABOUT]->(entity)
RETURN utterance.text, score, ep.title,
  collect(DISTINCT {claim: claim.text, entity: entity.name})
ORDER BY score DESC
\end{lstlisting}

\subsubsection{Text-to-Cypher with Self-Healing}

For queries that don't fit templates, the system generates Cypher from natural language with full graph schema, few-shot examples, and read-only constraints. When a generated query fails, the error and schema are fed back to the LLM for self-correction---up to 2 retries.

\subsubsection{Reciprocal Rank Fusion}

Results from multiple strategies are fused via RRF~\cite{cormack2009rrf}:
\begin{equation}
    \text{RRF\_score}(d) = \sum_{r \in R} \frac{1}{k + \text{rank}_r(d)}, \quad k = 60
\end{equation}
This combines rankings from different strategies without requiring score calibration.

\subsection{MCP Server}

The server exposes 6 tools: \texttt{search\_lectures} (semantic search + fan-out), \texttt{find\_claims} (filtered claim search), \texttt{compare\_entities} (side-by-side claims), \texttt{topic\_timeline} (chronological evolution), \texttt{run\_cypher} (NL $\rightarrow$ Cypher for ad-hoc queries), and \texttt{get\_schema} (schema self-discovery). Design principles: domain tools for common queries (optimized, fast, easy to use correctly); generic tools for exploration (flexibility for unforeseen queries); capped results (prevent context overflow); structured data returned (client LLM handles generation).

\subsection{Cost and Optimizations}

Three mechanisms reduce costs: (1) \textbf{Prompt caching} (50\% discount)---system prompt + schema identical across chunks, sequential processing keeps prefix cached; (2) \textbf{Tiered model usage}---GPT-4o-mini handles 80\% of calls at $\sim$10$\times$ lower cost; (3) \textbf{Local caching}---hash of each call stored on disk, reprocessing costs \$0.

Estimated cost per episode ($\sim$1 hour, $\sim$17K tokens): Passes 1--2 $\sim$\$0.02; Pass 3 $\sim$\$0.04; Pass 4 $\sim$\$0.02; Embeddings $\sim$\$0.01; Community summarization (amortized) $\sim$\$0.01. \textbf{Total per episode: $\sim$\$0.10. For 134 episodes: $\sim$\$13--15}---orders of magnitude below Microsoft GraphRAG's estimated \$33K at similar scale.


% ============================================================
% REFERENCES
% ============================================================
\bibliographystyle{plainnat}
\bibliography{references}

\end{document}
